\documentclass[11pt,a4paper]{article}
\usepackage[T1]{fontenc}
\usepackage[utf8]{inputenc}
\usepackage{lmodern}
\usepackage{amsmath,amssymb,amsthm}
\usepackage{xcolor}
\usepackage[margin=27mm]{geometry}
\usepackage[colorlinks=true,urlcolor=blue,citecolor=blue]{hyperref}
\newtheorem{theorem}{Theorem}
\newtheorem{proposition}[theorem]{Proposition}
\newtheorem{corollary}[theorem]{Corollary}
\theoremstyle{definition}
\newtheorem{definition}[theorem]{Definition}
\newtheorem{property}[theorem]{Property}
\theoremstyle{plain}
\newsavebox{\statementbox}
\newenvironment{statementframe}{%
  \par\addvspace{0.55\baselineskip}\noindent
  \begingroup\setlength{\fboxrule}{0.3pt}\setlength{\fboxsep}{7pt}%
  \begin{lrbox}{\statementbox}%
  \begin{minipage}{\dimexpr\linewidth-2\fboxsep-2\fboxrule\relax}%
}{%
  \end{minipage}\end{lrbox}%
  {\color{black!45}\fbox{\color{black}\usebox{\statementbox}}}%
  \endgroup\par\addvspace{0.45\baselineskip}%
}
\DeclareMathOperator{\Tr}{Tr}
\newcommand{\ket}[1]{\lvert#1\rangle}
\newcommand{\bra}[1]{\langle#1\rvert}
\newcommand{\norm}[1]{\left\lVert#1\right\rVert}
\setlength{\emergencystretch}{2em}
\setlength{\parindent}{0pt}
\setlength{\parskip}{0.3em plus 0.1em minus 0.1em}
\title{Exact ensemble-steering defect under partial environmental access}
\author{Fable \& Astra\\[0.35em]
\normalsize Research led by Jean-Philippe Bruneton\\[-0.1em]
\small Universit\'e Paris Cit\'e, France}
\date{17 November 2026; revised 19 November 2026}

\begin{document}
\maketitle
\begin{abstract}
Quantum steering allows an observer to prepare different ensembles for a
distant system by choosing a measurement on a system that purifies it. This
control is complete when the observer holds the whole purifier, but it can
fail when part of the purifying information has escaped into an inaccessible
environment. We determine exactly how large this failure is for a qubit state
with two branches. If the branches have weights $p$ and $q$ and their hidden
environmental records have overlap $\kappa$, the worst-case
ensemble-preparation error is $(1-\kappa)\sqrt{pq}$. Thus the lost steering
power is precisely the branch coherence recorded outside the observer's
reach. A binary target ensemble and a two-outcome measurement attain this
minimax value; when the hidden records are orthogonal, the value reduces to
half the known $\ell_1$ coherence of assistance of a qubit. For a chain of
summable controlled rotations, the same formula gives matching spatial bounds
and an asymptotically sharp geometric decay rate. The state is also the unique
ground state of a gapped quasi-local parent Hamiltonian. The proofs are short
and use standard tools; the contribution is the exact value of this
operational defect in an explicit model, and no general quantitative theory
of partial steering is claimed.
\end{abstract}

\section{Introduction}
Quantum steering is the ability to control remotely how a mixed state is
prepared. For example, suppose that Alice and Bob share the Bell state
$(\ket{00}+\ket{11})/\sqrt{2}$. If Bob measures his qubit in the computational
basis, Alice receives either $\ket0$ or $\ket1$. If he measures in the
diagonal basis, she instead receives either
$\ket+=(\ket0+\ket1)/\sqrt{2}$ or
$\ket-=(\ket0-\ket1)/\sqrt{2}$. Alice's average state is $I/2$ in both cases;
Bob's choice changes its ensemble decomposition, not the average state itself.

This freedom extends beyond the Bell pair. A mixed state admits many ensemble
decompositions, and the Schr\"odinger--HJW theorem (after Schr\"odinger,
and Hughston, Jozsa and Wootters) shows that an observer holding its entire
purifying system can prepare every finite decomposition into pure states by
measuring the purifier in a suitable basis, after enlarging it if
necessary \cite{Kirkpatrick}. Decompositions into mixed states are
coarse-grainings of pure ones and are prepared by merging outcomes. Perfect
steering therefore relies on access to all the information that purifies the
state.

The question studied here is what remains possible when that access is only
partial. We split the purifier into a subsystem $B$ that Bob can measure and a
hidden environment $C$. If $C$ retains a record of the alternative branches,
some ensembles of Alice's state can no longer be prepared from $B$. We measure
the resulting loss by the largest, over all target ensembles, of the smallest
trace-distance error achievable by a measurement on $B$.

For a qubit state with two branches, this worst-case error has a simple exact
meaning. Let $p$ and $q$ be the branch weights and let $\kappa$ be the overlap
between their records in $C$. The steering defect is
$(1-\kappa)\sqrt{pq}$: it vanishes when the hidden environment carries no
branch information and grows with the coherence that has become inaccessible.
The minimax formulation below makes this statement precise and identifies a
binary target ensemble and a two-outcome measurement that attain the value.
The same parameters acquire a spatial meaning when progressively larger
portions of an environment become accessible.

The exact formula has two consequences. First, it relates the steering defect
to a natural trace-distance defect for state-preserving conditional
expectations. Second, it turns the tail of a summable controlled-rotation
chain into explicit upper and lower bounds on the benefit of increasing
environmental access. The resulting decay is quadratic in the small tail
angles, although the unitary truncation error is linear.

For a state $\rho_{AB}$ with marginal $\rho_A$, a target ensemble is a finite
family of positive operators $\{\lambda_i\}$ satisfying
$\sum_i\lambda_i=\rho_A$; the probabilities are included in the traces of the
$\lambda_i$. We define
\begin{equation}\label{defect}
s_B(\rho)=\sup_{\{\lambda_i\}:\,\sum_i\lambda_i=\rho_A}
\inf_{\{E_i\}\,\mathrm{POVM\ on}\ B}
\frac12\sum_i\norm{\lambda_i-\Tr_B[(I\otimes E_i)\rho_{AB}]}_1.
\end{equation}
The target and measurement outcomes carry the same labels, and the
positive-operator-valued measure (POVM) may depend on the target ensemble.
This quantity is distinct from robustness against local-hidden-state models:
it asks how well one measurement system can realize every decomposition of a
fixed marginal.

\section{Related work}
Schr\"odinger's steering observation was given an operational definition and
placed in a strict hierarchy between entanglement and Bell nonlocality by
Wiseman, Jones and Doherty \cite{Wiseman}. Their question is whether the
conditional states admit a local-hidden-state model. The minimax
\eqref{defect} asks a different, preparation-oriented question: how closely
one fixed measurement system can realize every decomposition of a fixed
marginal. Its starting point is the Schr\"odinger--HJW theorem, whose history
and exact statement are reviewed by Kirkpatrick \cite{Kirkpatrick}; the
defect measures how this theorem fails once part of the purifier is removed.

For two-qubit states, the steering ellipsoid is the set of Bloch vectors to
which measurements on one qubit can steer the other \cite{Jevtic}. Jevtic,
Pusey, Jennings and Rudolph call steering complete when every convex
decomposition of the marginal Bloch vector into points of this ellipsoid is
realizable, and they show that this holds for every non-degenerate ellipsoid,
hence for every entangled two-qubit state \cite[Appendix~C]{Jevtic}. The
state $\rho_{AB}$ studied below is entangled for $\kappa>0$ and is therefore
completely steerable in their sense. Equation \eqref{defect} also counts
decompositions whose members lie outside the ellipsoid, and its exact value
quantifies how far these targets are from every realizable ensemble.

The coherence of assistance supplies the extremal ensemble. Zhao, Ma, Quan,
Fan and Pereira define the $\ell_1$ coherence of assistance $C_a^{\ell_1}$ as
the largest average $\ell_1$ coherence over pure-state decompositions, prove
$C_a^{\ell_1}(\rho)\leq\sum_{i\neq j}\sqrt{\rho_{ii}\rho_{jj}}$ by the
Cauchy--Schwarz inequality, and show that equality holds for qubits, with the
binary ensemble $\sqrt p\ket0\pm\sqrt q\ket1$ optimal
\cite[Theorem~1, Corollary~1, Lemma~1]{Zhao}. When the hidden records are
orthogonal ($\kappa=0$), every state Bob prepares is diagonal, and
Theorem~\ref{exact} reduces to $s_B=\sqrt{pq}=C_a^{\ell_1}(\rho_A)/2$. The
new ingredient for $\kappa>0$ is the matching lower bound over all
measurements on $B$.

Branch states whose decoherence factors are products of overlaps over
subenvironments are standard in quantum Darwinism, where they serve to
quantify the redundancy of records in environment fragments
\cite[Secs.~III--IV]{BlumeKohout} and \cite{Zurek}. The spatial model of
Section~\ref{sec:spatial} uses the same structure for a different task: it
measures the steering power gained as larger portions of the environment
become accessible.

We did not find the minimax value of Theorem~\ref{exact}, the comparison
with $\delta_*$ in Section~\ref{sec:ce}, or the spatial rate of
Corollary~\ref{spatial} in the literature consulted. Their proofs combine
the standard tools recalled above.

\section{Exact minimax result}
Consider the pure state
\begin{equation}\label{state}
\ket\Psi=\sqrt p\ket0_A\ket{b_0}_B\ket u_C
+\sqrt q\ket1_A\ket{b_1}_B\ket v_C,
\qquad p+q=1,\quad 0<q\leq p,
\end{equation}
where $b_0,b_1$ are orthonormal and $u,v$ are unit vectors. The system $C$ is
inaccessible. A phase choice gives $\langle v|u\rangle=\kappa\in[0,1]$.
On the logical support of $B$,
\[
\rho_{AB}=p\ket{00}\bra{00}+q\ket{11}\bra{11}
+\kappa\sqrt{pq}(\ket{00}\bra{11}+\ket{11}\bra{00}).
\]
Compressing a physical POVM to this support gives a qubit POVM. Conversely,
every logical POVM extends to the physical space by assigning the orthogonal
complement to one outcome.

\begin{statementframe}
\begin{theorem}\label{exact}
For the state in \eqref{state},
\[
\boxed{s_B(\rho)=(1-\kappa)\sqrt{pq}.}
\]
The supremum is attained by
$\lambda_\pm=\ket{\phi_\pm}\bra{\phi_\pm}/2$, where
$\ket{\phi_\pm}=\sqrt p\ket0\pm\sqrt q\ket1$.
\end{theorem}
\end{statementframe}
\begin{proof}
Set $D=\operatorname{diag}(\sqrt p,\sqrt q)$ and write
$\lambda_i=\begin{pmatrix}x_i&z_i\\\bar z_i&y_i\end{pmatrix}$.
The logical POVM $E_i=(D^{-1}\lambda_iD^{-1})^T$ is positive and sums to
$I$. It prepares
\[
\mu_i=\begin{pmatrix}x_i&\kappa z_i\\
\kappa\bar z_i&y_i\end{pmatrix}.
\]
Its cost is $(1-\kappa)\sum_i|z_i|$. Positivity and the Cauchy--Schwarz
inequality give
$\sum_i|z_i|\leq\sum_i\sqrt{x_i y_i}\leq\sqrt{pq}$, as in the qubit case of
\cite[Theorem~1]{Zhao}, establishing an upper bound for every finite target
ensemble.

For the matching lower bound, use the binary target stated in the theorem.
Every effect $0\leq E\leq I$ satisfies $|E_{01}|\leq1/2$, because
$\norm{2E-I}\leq1$. For $\Delta=\lambda_+-\mu_+$,
\[
\norm\Delta_1\geq\Tr(X\Delta)
=\sqrt{pq}-2\kappa\sqrt{pq}\,\Re E_{10}
\geq(1-\kappa)\sqrt{pq},
\qquad X=\begin{pmatrix}0&1\\1&0\end{pmatrix}.
\]
The two differences are opposite because both ensembles sum to $\rho_A$;
their total cost is therefore $\norm\Delta_1$. The effects
$E_\pm=(I\pm X)/2$ attain equality.
\end{proof}

\section{Conditional-expectation defect}\label{sec:ce}
When Bob holds the whole purifier of $\rho_{AC}$, perfect steering of $A$
from $B$ would follow if $\rho_{AC}$ were a product $\rho_A\otimes\tau$, that
is, if the conditional expectation
$E_\tau(X\otimes Y)=X\Tr(\tau Y)\otimes I$ of the matrix algebra of $AC$
onto that of $A$ preserved $\rho_{AC}$. Every conditional expectation of a
full matrix algebra $M_A\otimes M_C$ onto $M_A\otimes I$ has this form for
some state $\tau$. Define the optimal state-preservation defect
$\delta_*=\inf_\tau\frac12\norm{\rho_{AC}-\rho_A\otimes\tau}_1$. Here
\[
\rho_{AC}=p\ket0\bra0\otimes\ket u\bra u
+q\ket1\bra1\otimes\ket v\bra v.
\]
Writing $U=\ket u\bra u$ and $V=\ket v\bra v$ gives
\begin{align*}
2\delta_\tau&=p\norm{U-\tau}_1+q\norm{V-\tau}_1\\
&\geq q\norm{U-V}_1+(p-q)\norm{U-\tau}_1.
\end{align*}
The infimum is attained at $\tau=U$, and hence
\begin{equation}\label{delta}
\boxed{\delta_*=q\sqrt{1-\kappa^2},\qquad
s_B=\sqrt{pq}\left(1-\sqrt{1-(\delta_*/q)^2}\right).}
\end{equation}
Faithful states approximating $U$ preserve the same infimum if faithfulness
is required. The second equality assumes fixed $q$ and
$0\leq\delta_*\leq q$.

At $\kappa=0$ with $q\to0$, one has
$s_B=\sqrt{\delta_*(1-\delta_*)}$, which rules out a uniform linear upper
bound in $\delta_*$. For fixed $q>0$ and $\kappa\to1$,
\[
s_B\sim\frac{\sqrt{pq}}{2q^2}\delta_*^2.
\]
The two limits identify different scaling regimes: a general square-root
bound can be far from sharp when leakage is small and the Schmidt weights
remain fixed.

\section{Spatially increasing access}\label{sec:spatial}
Place a qubit $A$, a flag qubit $B_0$, and qubits $E_j$ for $j\geq1$ in the
state obtained from
$(\sqrt p\ket{00}+\sqrt q\ket{11})_{AB_0}\otimes_j\ket0_{E_j}$ by the
controlled rotations
\[
U_j=\ket0\bra0_A\otimes I+\ket1\bra1_A\otimes R(\theta_j),
\qquad
R(\theta)=\begin{pmatrix}\cos\theta&-\sin\theta\\
\sin\theta&\cos\theta\end{pmatrix}.
\]
Assume $0\leq\theta_j<\pi/2$ and $\sum_j\theta_j<\infty$. The commuting
product $U=\prod_jU_j$ converges in norm in the quasi-local algebra because
$\norm{U_j-I}=2\sin(\theta_j/2)\leq\theta_j$. Bob controls
$B_R=B_0E_1\cdots E_R$, and the remaining qubits form $C_R$.

\begin{statementframe}
\begin{corollary}\label{spatial}
For $V_R=\sum_{j>R}\theta_j^2$,
\begin{align}
s_{B_R}&=\sqrt{pq}\left(1-\prod_{j>R}\cos\theta_j\right),
\label{spatialexact}\\
\sqrt{pq}(1-e^{-V_R/2})&\leq s_{B_R}\leq\frac{\sqrt{pq}}2V_R.
\label{bounds}
\end{align}
Consequently, $s_{B_R}\sim\sqrt{pq}V_R/2$ whenever $V_R\to0$ through
positive values. A vanishing tail gives zero defect exactly.
\end{corollary}
\end{statementframe}
\begin{proof}
The inaccessible branches have overlap $\prod_{j>R}\cos\theta_j$, so
Theorem~\ref{exact} gives \eqref{spatialexact}. For a finite sequence,
$1-\prod_j a_j\leq\sum_j(1-a_j)$ when $0\leq a_j\leq1$. The inequalities
$1-\cos\theta\leq\theta^2/2$ and
$\cos\theta\leq e^{-\theta^2/2}$ on $[0,\pi/2)$ then give
\eqref{bounds}; convergence permits passage to the infinite product. The
ratio of the two bounds tends to one as $V_R\to0$.
\end{proof}

For $\theta_j=\theta_0r^j$ with $0<r<1$ and $0<\theta_0<\pi/2$,
\begin{equation}\label{rate}
s_{B_R}\sim\frac{\sqrt{pq}\theta_0^2}{2(1-r^2)}r^{2(R+1)}.
\end{equation}
The unitary truncation obeys
$\norm{U-U_{\leq R}}\leq\theta_0r^{R+1}/(1-r)$, so the ensemble task
decays to second order while the operator-norm truncation error is first
order. The sufficient access radius for $s_{B_R}\leq\varepsilon$ is
\[
R\geq\max\left\{0,\left\lceil
\frac{\log[\sqrt{pq}\theta_0^2/(2(1-r^2)\varepsilon)]}
{2|\log r|}-1\right\rceil\right\}.
\]
The optimal binary measurement acts on the span of the two branches of
$B_R$. It can be implemented by undoing the accessible controlled rotations
and measuring $B_0$ in the eigenbasis of $X$.

\clearpage
\section{A gapped quasi-local parent model}
On the product representation of the environmental qubits, define
\[
H_0=(I-\ket{\psi_p}\bra{\psi_p})_{AB_0}
+\sum_{j\geq1}\ket1\bra1_{E_j},
\qquad \ket{\psi_p}=\sqrt p\ket{00}+\sqrt q\ket{11}.
\]
This positive excitation-number operator has a unique ground state and gap
one. The unitarily equivalent operator $H=UH_0U^*$ on $U\mathcal D(H_0)$
has the same spectrum and the state above as its ground state. Each term
$U\ket1\bra1_{E_j}U^*-\ket1\bra1_{E_j}$ is supported on $\{A,E_j\}$ and
has norm at most $2\theta_j$. The transformed $AB_0$ term is quasi-local,
with truncation error at most $2\sum_{j>R}\theta_j$. Thus the parent model is
well defined for a summable angle sequence.

The interaction is inhomogeneous and long ranged, although summable. Its
ground state is related to a block-product state by a quasi-local unitary,
and no topological order is established. The partition $A,B_R,C_R$ keeps
Alice's system fixed while Bob's region grows, so the decay rate describes
the gain from enlarging the controlled region, not the dependence of
steering on the distance between two fixed observers. The model therefore
supplies an exactly solvable example, while a general quantitative theory of
steering under partial access, for instance in lattice systems with
topological order, remains open.

\section{Verification and scope}
The public
\href{https://github.com/jpbruneton/open_quantum_problems/blob/main/research/entries/steering-localization/checks/verify51.py}{Python checker}
verifies the partial-trace preparations, the binary witnesses, rational test
cases, spatial tails, and spectra of finite parent Hamiltonians; its
\href{https://github.com/jpbruneton/open_quantum_problems/blob/main/research/entries/steering-localization/data/verification51.json}{recorded output}
is included in the same public entry. These computations test finite
instances and complement the analytic proofs, which alone establish the
universal statements. Bibliographic priority is a separate question: the
search reported in Section~2 found the ingredients of the proofs but not the
exact minimax value, which remains open to independent assessment.

\clearpage
\begingroup\small
\begin{thebibliography}{99}
\setlength{\itemsep}{2pt}
\bibitem{Kirkpatrick}
K. A. Kirkpatrick, \emph{The Schr\"odinger--HJW theorem},
Found. Phys. Lett. \textbf{19}, 95--102 (2006);
\href{https://doi.org/10.1007/s10702-006-1852-1}{doi:10.1007/s10702-006-1852-1};
\href{https://arxiv.org/abs/quant-ph/0305068v3}{arXiv:quant-ph/0305068v3}.

\bibitem{Wiseman}
H. M. Wiseman, S. J. Jones and A. C. Doherty,
\emph{Steering, Entanglement, Nonlocality, and the Einstein--Podolsky--Rosen Paradox},
Phys. Rev. Lett. \textbf{98}, 140402 (2007);
\href{https://doi.org/10.1103/PhysRevLett.98.140402}{doi:10.1103/PhysRevLett.98.140402};
\href{https://arxiv.org/abs/quant-ph/0612147v3}{arXiv:quant-ph/0612147v3}.

\bibitem{Jevtic}
S. Jevtic, M. F. Pusey, D. Jennings and T. Rudolph,
\emph{Quantum Steering Ellipsoids}, Phys. Rev. Lett. \textbf{113}, 020402 (2014);
\href{https://doi.org/10.1103/PhysRevLett.113.020402}{doi:10.1103/PhysRevLett.113.020402};
\href{https://arxiv.org/abs/1303.4724v2}{arXiv:1303.4724v2}.

\bibitem{Zhao}
M.-J. Zhao, T. Ma, Q. Quan, H. Fan and R. Pereira,
\emph{$l_1$-norm coherence of assistance},
Phys. Rev. A \textbf{100}, 012315 (2019);
\href{https://doi.org/10.1103/PhysRevA.100.012315}{doi:10.1103/PhysRevA.100.012315};
\href{https://arxiv.org/abs/2002.03823}{arXiv:2002.03823}.

\bibitem{BlumeKohout}
R. Blume-Kohout and W. H. Zurek,
\emph{Quantum Darwinism: Entanglement, branches, and the emergent classicality
of redundantly stored quantum information},
Phys. Rev. A \textbf{73}, 062310 (2006);
\href{https://doi.org/10.1103/PhysRevA.73.062310}{doi:10.1103/PhysRevA.73.062310};
\href{https://arxiv.org/abs/quant-ph/0505031}{arXiv:quant-ph/0505031}.

\bibitem{Zurek}
W. H. Zurek, \emph{Quantum Darwinism}, Nat. Phys. \textbf{5}, 181--188 (2009);
\href{https://doi.org/10.1038/nphys1202}{doi:10.1038/nphys1202};
\href{https://arxiv.org/abs/0903.5082}{arXiv:0903.5082}.
\end{thebibliography}
\endgroup
\end{document}
